\chapter{Compact Riemann Surfaces and Genus}
\label{ch:riemann-surfaces}

A Riemann surface is the natural setting for one-variable complex analysis after
we stop insisting that all functions live on open subsets of \(\C\).  It is a
topological surface with charts to \(\C\) whose transition functions are
holomorphic.  Compact Riemann surfaces are the complex-analytic analogues of
closed algebraic curves.

\begin{definition}
A compact Riemann surface \(X\) has genus \(g\) if its underlying oriented
topological surface has \(g\) handles.  Equivalently,
\(\chi(X)=2-2g\).
\end{definition}

The midterm computation that the real torus has Euler characteristic \(0\) is
therefore already telling us that a complex torus has genus \(1\).  This number is
not cosmetic.  It is exactly the number that makes Riemann--Roch read
\(\ell(D)=\deg D\) for positive divisors on an elliptic curve.

\begin{proposition}
\label{prop:compact-functions}
Every holomorphic function on a compact connected Riemann surface is constant.
\end{proposition}

\begin{proof}
The image of a compact surface under a holomorphic function is compact in \(\C\),
so the function has a maximum absolute value.  The maximum modulus principle then
forces the function to be locally constant, hence constant.
\end{proof}

\begin{definition}
A meromorphic function on a Riemann surface is a holomorphic map to
\(\CP^1\).  Its poles are the points mapping to \(\infty\).
\end{definition}

\begin{keyidea}{Compactness changes the rules}
On a compact surface, nonconstant holomorphic functions disappear, but
meromorphic functions become structured maps to \(\CP^1\).  The map
\(\wp:T\to\CP^1\) is exactly such a map.
\end{keyidea}

\section*{Exercises}
\addcontentsline{toc}{section}{Exercises}

\exerciseintent{Convert Euler characteristic into genus in the genus-one case.}
\begin{exercise}
Substitute \(\chi=0\) into \(\chi=2-2g\) and solve the resulting linear equation
for \(g\).
\end{exercise}

\exerciseintent{Identify exactly where compactness enters the maximum modulus argument.}
\begin{exercise}
Use the concrete function \(f(z)=z\) on \(\C\).  Compute why \(|f|\) has no
maximum on \(\C\), and use this calculation to identify the step in the maximum
modulus proof that fails on noncompact surfaces.
\end{exercise}

\exerciseintent{Use compactness and local normal forms to prove zeros and poles are finite.}
\begin{exercise}
Suppose a meromorphic function on a compact surface had infinitely many zeros or
poles.  Use compactness to produce an accumulation point, then use the local
forms \((z-p)^m u(z)\) and \((z-p)^{-m}u(z)\) to compute why this contradicts
isolation of zeros and poles.
\end{exercise}

\exerciseintent{Recover the torus genus by a concrete cell count on the glued parallelogram.}
\begin{exercise}
The quotient of a parallelogram by identifying opposite sides is a topological
torus.  Count the cell classes \(F,E,V\) after the identifications, compute
\[
  \chi=F-E+V,
\]
and then use \(\chi=2-2g\) to compute the genus.
\end{exercise}
